\chapter{The Shuffle}
\label{ch:shuffle}

\begin{chapterintro}
Every \texttt{JOIN}, \texttt{GROUP BY}, \texttt{DISTINCT}, window function, and
repartition puts a wall between two stages: the shuffle. It writes every
partition to disk, moves it across the network, and re-reads it. It is where most
of a query's time goes. This chapter walks the write path and the read path,
shows how to size shuffle partitions from data volume instead of accepting the
default 200, and how to find and remove exchanges the query does not need.
\end{chapterintro}

\section{What a partition is: input partitions versus shuffle partitions}
\tbd

\section{The shuffle write path}
\tbd

\section{The shuffle read path}
\tbd

\section{Sizing shuffle partitions}
\tbd

\section{Shuffle compression}
\tbd

\section{Push-based shuffle}
\tbd

\section{Finding and removing unnecessary exchanges}
\tbd

\section{Summary}
\tbd

\exercisehead
\begin{exercises}
\item From a stage's shuffle-read size in the UI, compute a target partition
      count for a 128~MB task size, then verify against the observed task sizes.
\end{exercises}
